\section{元の変数を新しい変数で表す}\label{節:逆向きの置換}

合成形がそのまま見えない場合でも，根号などを新しい変数にしてから，元の変数を解き直すと式が簡単になることがある。このときは，$t=g(x)$ から $x$ を $t$ の関数として表せる範囲を選ぶ。

\begin{example}{根号を新しい変数にする}
$x>0$ において，$\int\frac{\sqrt{x}}{1+\sqrt{x}}\dd x$ を求めよ。
\end{example}

$t=\sqrt{x}$ と置けば，$t>0$ で $x=t^2$，$\dd x=2t\dd t$ となる。分子だけでなく分母と $\dd x$ も書き換えて
\[
 \int\frac{\sqrt{x}}{1+\sqrt{x}}\dd x
 =\int\frac{t}{1+t}\,2t\dd t
 =2\int\frac{t^2}{t+1}\dd t.
\]
根号はなくなったが分数が残った。ここで前章の割り算を使い，$t^2=(t+1)(t-1)+1$ と分ける。
\[
\begin{aligned}
 2\int\frac{t^2}{t+1}\dd t
 &=2\int\left(t-1+\frac1{t+1}\right)\dd t\\
 &=t^2-2t+2\log(t+1)+C.
\end{aligned}
\]
$t+1>0$ なので，ここでは対数の絶対値を外せる。$t=\sqrt{x}$ を戻すと，答えは
\[
 x-2\sqrt{x}+2\log(1+\sqrt{x})+C.
\]
微分して確かめると
\[
 1-\frac1{\sqrt{x}}+\frac1{(1+\sqrt{x})\sqrt{x}}
 =1-\frac1{1+\sqrt{x}}
 =\frac{\sqrt{x}}{1+\sqrt{x}}
\]
となる。置換が積分を一度で終わらせなくても，根号を有理式へ変えるという目的は達している。

この例で $x=t^2$ と書けたのは，$t=\sqrt{x}>0$ と決めていたからである。逆に $t=x^2$ から $x$ を求めるなら，$x=\sqrt{t}$ と $x=-\sqrt{t}$ の二つがある。$x>0$ の区間では前者，$x<0$ の区間では後者を使う。このように，逆関数を一つに決めるために選ぶ範囲を\textbf{枝}と呼ぶことがある。

区間上の連続関数で，逆関数を持つことと狭義単調であることは同値である。ここで狭義単調とは，異なる入力に対して出力が必ず一方向に増えるか，必ず一方向に減ることをいう。したがって，$t=g(x)$ から一つの関数 $x=h(t)$ を使って戻したいなら，そのような区間を選ばなければならない。さらに逆関数の微分を $h'(t)=\frac1{g'(h(t))}$ と書く場合には，分母が $0$ でないことも必要である。

これに対し，\sectionref{節:合成形の置換}の公式では逆関数を作っていない。$F(g(x))$ をそのまま微分したので，単調性を仮定しなかった。定積分の公式
\[
 \int_\alpha^\beta f(h(t))h'(t)\dd t
 =\int_{h(\alpha)}^{h(\beta)}f(x)\dd x
\]
自体も，同じ合成形の公式であり，$h$ の単調性を必要としない。単調性が必要になるのは，逆関数を一つ選んで式を書き直す段階である。この二つを区別すれば，置換に単調性が要る場合と要らない場合を混同せずに済む。

